\section*{Chapter \ref{ch:m1:index-and-single-stock-futures-worldwide} --- Equity Index and Single-Stock Futures Worldwide}

\begin{solution}{exo:m1:index-and-single-stock-futures-worldwide:1}
\$300\,000, \euro54\,000 and HK\$1\,300\,000; a 1\% move is \$3\,000, \euro540 and
HK\$13\,000 per contract.
\end{solution}

\begin{solution}{exo:m1:index-and-single-stock-futures-worldwide:2}
$5\,400 \times (1 + 0.025 \times 0.25) - 5\,412.0 = 5\,433.75 - 5\,412.0 = 21.75$
points.
\end{solution}

\begin{solution}{exo:m1:index-and-single-stock-futures-worldwide:3}
$(158.2 - 171.5) \times \text{\euro}100 \times 50 = -\text{\euro}66\,500$.
\end{solution}

\begin{solution}{exo:m1:index-and-single-stock-futures-worldwide:4}
$+2.08\%$, $-1.46\%$, $-1.78\%$, $-0.90\%$, $-0.61\%$; over five years
$(163.5/168.0)^{1/5} - 1 = -0.54\%$ a year. The strip is a price, not a
forecast: it is what the sellers of dividend risk, mostly hedgers of
structured products, accept to be rid of it. The gap to the analyst's
forecast is the premium a buyer earns for carrying that risk, together with
the possibility that the analyst is wrong.
\end{solution}

\begin{solution}{exo:m1:index-and-single-stock-futures-worldwide:5}
Hong Kong's three sessions lie inside the hours of the other two: $2.75 + 3.5
+ 10 = 16.25$ hours with all three open. Nothing is open from 21:00 to 22:00
UTC, the American contract's daily break. In December Frankfurt and Chicago
are one hour later in UTC (their clocks go back; UTC and Hong Kong do not
move): the European bar runs from 01:10 to 21:00 and the break falls at
22:00.
\end{solution}

\begin{solution}{exo:m1:index-and-single-stock-futures-worldwide:6}
$200 \times \text{\euro}10 \times 5\,400 \times 0.0015 \times 2 =
+\text{\euro}32\,400$. The long pays the spread as financing; having locked 40
when the market now charges 55, the position is worth the difference for
the two remaining years.
\end{solution}

\begin{solution}{exo:m1:index-and-single-stock-futures-worldwide:7}
On 16 December the empty hour is 22:00 to 23:00 UTC, against 21:00 to 22:00
in September. During the week of 14 September the set changes at each
session boundary: Sunday 22:00 UTC (the American contract opens), then each
day 00:10 (Europe opens), 01:15, 04:00, 05:00, 08:30, 09:00 and 19:00 (Hong
Kong's three sessions), 20:00 (Europe closes), 21:00 and 22:00 (the American
break), until Friday 21:00, when the last contract closes.
\end{solution}

\begin{solution}{exo:m1:index-and-single-stock-futures-worldwide:8}
The statistic counts contracts, and contracts differ in size by factors of a
thousand. The 2025 total fell 42\% while futures volume rose 8.6\%: the count
is dominated by very small equity options in a few markets, and moved with
regulatory changes to those contracts, not with the amount of risk
transferred. A firm should rank markets by notional or, better, by the
revenue available to its kind of strategy: volume in money, spreads, fees,
access and competition. Some of the markets behind the 63\% are also closed,
or nearly so, to foreign proprietary firms.
\end{solution}

\begin{solution}{pb:m1:index-and-single-stock-futures-worldwide:1}
\textbf{1.} $171.5/5\,400 = 3.18\%$.
\textbf{2.} $+2.08\%$.
\textbf{3.} $-0.54\%$ a year.
\textbf{4.} $168.0 \times 1.04^5 = 204.4$; the future at 163.5 is 20\% below.
\textbf{5.} Dealers hedging the long-dividend exposure that structured
products leave them with; they sell at a discount because nobody else is a
natural seller's counterpart in size for a date five years away.
\textbf{6.} $200 \times 163.5 \times 100 = \text{\euro}3.27$ million.
\textbf{7.} $(204.4 - 163.5) \times 20\,000 = +\text{\euro}818\,000$.
\textbf{8.} $168.0 \times 0.7 = 117.6$; $(117.6 - 163.5) \times 20\,000 =
-\text{\euro}918\,000$.
\textbf{9.} $(110 - 163.5) \times 20\,000 = -\text{\euro}1.07$ million, in
cash, within the month, on a position of \euro3.27 million. Being right in 2031
requires surviving the \omterm{def:m1:financing:margincall}{margin calls} of \cref{ch:m1:margin} until then.
\textbf{10.} $20\% \times 171.5 = 34.3$ points.
\textbf{11.} The future's fair value rises by 34.3 points: $500 \times 10
\times 34.3 = +\text{\euro}171\,500$. A long future is short dividends, so a
cut is a gain.
\textbf{12.} Buy 50 \omterm{def:m1:index-and-single-stock-futures-worldwide:divfut}{dividend futures} of 2027: $50 \times 100 \times 34.3 =
\text{\euro}171\,500$.
\textbf{13.} The products it sells pay investors the index's price return,
often with a guarantee; the desk hedges by holding the index, through shares
or futures, and collects the dividends while owing none. Its profit depends
on dividends for the product's life, five to ten years.
\textbf{14.} The \omterm{def:m1:index-and-single-stock-futures-worldwide:trf}{total return future} pays realised dividends instead of
embedding an estimate of them, so the hedge has no dividend exposure to
begin with, and it fixes the financing spread for the product's life.
\textbf{15.} No. It is the price at which dividend risk clears between
structural sellers and the capital willing to hold it, and moves with that
capital's appetite as much as with dividend expectations.
\textbf{16.} That the holders are leveraged or marked-to-market investors who
are forced out when prices fall and margins rise, so that the contract
behaves in a crisis like a leveraged equity position, not like the slow
variable it settles on.
\textbf{17.} Nothing tradable: the 2031 dividends do not exist yet. The
contract can be hedged only approximately, with the index, with other
points of the strip, or with options; its risk is largely unhedgeable, which
is why it carries a premium.
\textbf{18.} Less liquid. They exist for the same reason at the level of one
share: option and structured-product desks are long single-stock dividends
and need to sell them, and a dividend cut in one name is a large, discrete
risk.
\textbf{19.} \textbf{$+2.1\%$ for the first year; $-0.5\%$ a year over five
years.}
\textbf{20.} The dividend leg of owning equity, from its price and its
financing.
\end{solution}

\begin{solution}{iq:m1:index-and-single-stock-futures-worldwide:1}
The E-mini S\&P 500 in Chicago, the EURO STOXX 50 future in Frankfurt, the
Hang Seng or Nikkei futures in Asia. The E-mini expires on the third Friday
against an opening quotation; the EURO STOXX 50 future the same day at noon
against a ten-minute average; the Hang Seng future monthly, on the
second-last trading day, against an average of five-minute quotations.
\iqlookfor{one concrete expiry detail per contract.}
\end{solution}

\begin{solution}{iq:m1:index-and-single-stock-futures-worldwide:2}
A future settling on the dividends an index pays over a period, in \omterm{def:m1:index-and-single-stock-futures-worldwide:point}{index
points}. Sellers: dealers long dividends from structured products and
options books. Buyers: funds collecting the discount, and relative-value
traders against the index future, options and single-name dividends.
\iqlookfor{the structural seller.}
\end{solution}

\begin{solution}{iq:m1:index-and-single-stock-futures-worldwide:3}
If the dividend fell before my future's expiry, I gain: the future was
priced below spot-plus-interest by the expected dividend, the index will no
longer drop by that amount on the ex-date, so fair value rises by the
dividend in \omterm{def:m1:index-and-single-stock-futures-worldwide:point}{index points}. The spot index itself may fall on the news; that
is a separate, ordinary price move.
\iqlookfor{long future equals short dividends, separated from the price
reaction.}
\end{solution}

\begin{solution}{iq:m1:index-and-single-stock-futures-worldwide:4}
The term financing spread of equities, alone. The ordinary future bundles an
estimate of dividends and a financing rate to one expiry; the \omterm{def:m1:index-and-single-stock-futures-worldwide:trf}{total return
future} pays realised dividends and charges a realised overnight rate, is
quoted in the spread, and lists expiries out to ten years, so that a view or
a hedge on equity repo can be held as a cleared, margined position.
\iqlookfor{``quoted in the spread'' and the bank balance-sheet motive.}
\end{solution}

\begin{solution}{iq:m1:index-and-single-stock-futures-worldwide:5}
Sessions in exchange-local time with an IANA zone name, never as UTC
offsets; several sessions per product per day, allowed to cross midnight;
the weekdays on which each session opens; a holiday and half-day table per
exchange and year; expiry-day exceptions from the contract master; an
effective-date range on every row, since hours change. All queries answered
in UTC by one function that every system shares.
\iqlookfor{zone names not offsets; effective dating.}
\end{solution}

\begin{solution}{iq:m1:index-and-single-stock-futures-worldwide:6}
Convert contracts to notional with each contract's multiplier and price, and
for options to premium or delta-adjusted notional. Separate futures from
options. Then add what decides whether a firm can earn anything there:
typical spread in basis points, tick constraint, fee and tax per trade, who
is allowed to trade and through whom, and the share of volume that is
professional against retail. The ranking that results bears little
resemblance to the league table.
\iqlookfor{notional conversion first, then economics and access.}
\end{solution}
